\documentclass[
    language=english,
    doctype=public-comment,
    institution=none,
]{../../omnilatex}

\RequirePackage{config/document-settings}

\begin{document}

\maketitle

\section*{Public Comment on Proposed EPA Regulation}

\begin{description}
    \item[Regulatory Agency.] U.S. Environmental Protection Agency
    \item[Docket Number.] EPA-HQ-OPPT-2026-0042
    \item[Proposed Rule.] Standards for Per- and Polyfluoroalkyl Substances
          (PFAS) in Drinking Water
    \item[Comment Period.] April 15 -- June 15, 2026
    \item[Commenter.] Environmental Coalition for Clean Water
    \item[Date.] \today
\end{description}

\section{Background}

The Environmental Coalition for Clean Water (ECCW) submits these comments
in response to the EPA's proposed rule establishing maximum contaminant
levels (MCLs) for six PFAS compounds in public drinking water systems. We
support the EPA's effort to address PFAS contamination, which poses
significant risks to public health and the environment. However, we believe
the proposed standards should be strengthened in several key areas.

PFAS chemicals, often referred to as ``forever chemicals,'' persist in the
environment and accumulate in human tissue. Research has linked PFAS exposure
to cancer, thyroid disease, immune system suppression, and developmental
delays in children. The proposed MCLs of 4 parts per trillion (ppt) for
PFOA and PFOS represent a meaningful step forward, but emerging evidence
suggests that even lower thresholds may be necessary to protect public
health.

\section{Analysis of Proposed Standards}

\subsection{Adequacy of MCL Levels}

\begin{table}[h]
    \centering
    \begin{tabular}{lccc}
        \toprule
        \textbf{Substance} & \textbf{Proposed MCL} & \textbf{ECCW Recommendation} & \textbf{Scientific Basis} \\
        \midrule
        PFOA & 4 ppt & 1 ppt & ATSDR minimal risk level \\
        PFOS & 4 ppt & 2 ppt & CDC bioaccumulation data \\
        PFHxS & 10 ppt & 4 ppt & Emerging epidemiology \\
        PFNA & 10 ppt & 5 ppt & Animal study no-observed-adverse-effect \\
        HFPO-DA & 10 ppt & 5 ppt & Drinking water equivalency \\
        PFBS & 2,000 ppt & 1,000 ppt & Relative potency factor \\
        \bottomrule
    \end{tabular}
    \caption{Comparison of proposed MCLs and recommended levels.}
\end{table}

\subsection{Monitoring and Testing Requirements}

We support the requirement for quarterly monitoring of systems serving more
than 10,000 people. However, we recommend:

\begin{enumerate}
    \item \textbf{Expanded monitoring.} Systems serving 3,300 to 10,000
          people should be required to conduct semi-annual monitoring,
          rather than the proposed biennial schedule.
    \item \textbf{Point-of-entry testing.} Residential wells within 5
          miles of known contamination sites should be included in the
          monitoring program.
    \item \textbf{Method standardization.} The EPA should designate a single
          analytical method (e.g., EPA 537.1 modified) to ensure
          comparability of results across laboratories.
\end{enumerate}

\subsection{Implementation Timeline}

The proposed 5-year implementation timeline for large systems and 8-year
timeline for small systems is appropriate, provided adequate funding is
made available. We recommend:

\begin{itemize}
    \item The EPA prioritize a \$2 billion annual appropriation for the
          Drinking Water State Revolving Fund to support PFAS treatment
          infrastructure.
    \item Extended compliance deadlines be conditioned on demonstrated
          progress toward installation of treatment technology.
    \item Interim health advisories be issued immediately for communities
          with PFAS levels exceeding 10 ppt.
\end{itemize}

\section{Environmental Justice Considerations}

PFAS contamination disproportionately affects communities near military
bases, industrial facilities, and airports. We urge the EPA to:

\begin{enumerate}
    \item Map PFAS contamination sites against demographic data to
          identify environmental justice communities.
    \item Provide supplemental technical assistance to small and rural
          water systems that lack engineering capacity.
    \item Ensure that compliance costs are not passed through to low-income
          ratepayers without federal subsidy support.
\end{enumerate}

\section{Recommendations}

Based on our analysis, we respectfully submit the following recommendations:

\begin{enumerate}
    \item Lower MCLs for PFOA and PFOS to 1 ppt and 2 ppt respectively,
          consistent with the most protective health-based standards.
    \item Expand monitoring requirements to include mid-size systems and
          residential wells near contamination sites.
    \item Establish a dedicated federal fund of \$2 billion annually for
          PFAS remediation in drinking water systems.
    \item Issue immediate interim health advisories for communities with
          PFAS levels exceeding 10 ppt.
    \item Require annual reporting of PFAS levels in drinking water to
          state health departments for public disclosure.
\end{enumerate}

\section{Conclusion}

We commend the EPA for taking this critical step to protect public health
from PFAS contamination. The proposed rule establishes an important
regulatory framework, but the science demands stronger standards. We urge
the Agency to adopt the recommendations outlined in this comment to ensure
that all Americans have access to safe, clean drinking water.

\begin{description}
    \item[Submitted by.] Environmental Coalition for Clean Water
    \item[Contact.] Dr. Rachel Nguyen, Director of Policy
    \item[Phone.] (202) 555-0187
    \item[Email.] rnguyen@eccw.org
\end{description}

\end{document}
